\documentclass[11pt]{article}
\usepackage[a4paper,margin=22mm]{geometry}
\usepackage{fontspec}
\setmainfont{DejaVu Serif}
\setsansfont{DejaVu Sans}
\usepackage{microtype}
\usepackage{xcolor}
\usepackage{hyperref}
\usepackage{enumitem}
\usepackage{parskip}
\definecolor{Accent}{HTML}{971D32}
\definecolor{Muted}{HTML}{666666}
\hypersetup{colorlinks=true,urlcolor=Accent,linkcolor=Accent}
\setlist{leftmargin=1.6em,itemsep=0.3em,topsep=0.35em}
\setcounter{tocdepth}{2}
\renewcommand{\contentsname}{Contents}
\newcommand{\sectionrule}{\par\vspace{-0.5em}\noindent\textcolor{Accent}{\rule{\linewidth}{0.6pt}}\par}
\begin{document}

\begin{center}
{\sffamily\bfseries\color{Accent} TOEFL WORD OF THE DAY\par}
\vspace{8mm}
{\fontsize{42}{48}\selectfont\bfseries rigid\par}
\vspace{2mm}
{\Large\sffamily\color{Accent} /ˈrɪdʒɪd/\par}
\vspace{2mm}
{\small\sffamily Local audio phrase: rigid\par}
{\small\sffamily Audio file: audios/rigid.mp3\par}
\vspace{2mm}
{\large\sffamily\color{Muted} adjective\par}
\vspace{5mm}
{\sffamily stiff and unable to bend \textbullet\ strict or difficult to change\par}
\vspace{6mm}
{\large Rigid describes something that cannot bend physically or something such as a rule, system, or attitude that does not easily change.\par}
\vspace{6mm}
\begin{minipage}{0.84\textwidth}
\itshape “A rigid system may work well in normal conditions but fail when circumstances change.”
\end{minipage}\par
\vspace{10mm}
\textbf{Date:} August 15th 2026\quad\textbullet\quad \textbf{Level:} B1--mid-B2 TOEFL
\vfill
Created and compiled by \textbf{Amir Kooshky}\par
\vspace{2mm}
\href{https://t.me/KooshkyTOEFL}{Telegram Channel}\quad\textbullet\quad
\href{https://www.instagram.com/kooshkytoefl}{Instagram}\quad\textbullet\quad
\href{https://t.me/KooshkyTOEFL\_pv}{Personal Telegram}
\end{center}
\newpage
\tableofcontents
\newpage

\section{Meanings and usage}
\sectionrule
\subsection{Meaning 1: Stiff and unable to bend}
\textbf{Part of speech:} adjective

\textbf{IPA:} /ˈrɪdʒɪd/

\textbf{Pronunciation phrase:} rigid

\textbf{Local audio file:} audios/rigid.mp3

\textbf{Register:} neutral to formal; common in science, engineering, biology, medicine, and descriptions of materials

\textbf{Definition:} stiff and not able to bend, change shape, or move easily

\textbf{In simpler words:} hard to bend

\textbf{Typical pattern:} rigid + structure, frame, material, surface, or support

\subsubsection{Natural examples}
\begin{enumerate}
  \item The building is supported by a rigid steel frame.
  \item Unlike flexible cartilage, bone is relatively rigid.
  \item The container must be rigid enough to protect the equipment during transportation.
  \item The insect's rigid outer covering helps protect its internal organs.
  \item Engineers needed a lightweight but rigid material for the structure.
  \item The cold temperature made the plastic more rigid and easier to break.
  \item A rigid surface does not absorb impact as easily as a flexible one.
  \item The telescope requires a rigid support to prevent small movements from affecting the measurements.
\end{enumerate}

\subsubsection{Nuance and implication}
\textbf{Central nuance:} Rigid emphasizes resistance to bending or changing shape.

\textbf{Usual implication:} Rigidity can provide strength and stability, but it can also make an object more likely to crack or break under certain forces.

\textbf{Compared with flexible:} A flexible object can bend or change shape relatively easily. A rigid object strongly resists bending or deformation.

\subsubsection{Grammar notes}
\textbf{How to use it:} Use it before a noun, as in a rigid frame, or after a linking verb, as in The frame is rigid.

\textbf{What can follow it:} It commonly describes materials, structures, surfaces, containers, supports, and parts of the body.

\textbf{Position in a sentence:} Words such as highly, relatively, completely, and sufficiently can appear before rigid.

\subsubsection{Useful patterns}
\textbf{Pattern:} a rigid + material or structure

\textbf{Use:} Use this for something that strongly resists bending or changing shape.

\textbf{Example:} The device is enclosed in a rigid plastic shell.

\textbf{Pattern:} remain rigid

\textbf{Use:} Use this when something continues to keep its shape despite pressure or movement.

\textbf{Example:} The frame must remain rigid under heavy loads.

\textbf{Pattern:} rigid enough to + verb

\textbf{Use:} Use this when a certain level of stiffness is necessary for a purpose.

\textbf{Example:} The panel is rigid enough to support the weight.

\textbf{Pattern:} become more or less rigid

\textbf{Use:} Use this when the stiffness of a material changes.

\textbf{Example:} Some plastics become more rigid at low temperatures.

\subsubsection{Common wrong usage}
\paragraph{Correction 1}
\textbf{Wrong:} The metal is very rigidity.

\textbf{Correct:} The metal is very rigid.

\textbf{Why:} Use the adjective rigid after is. Rigidity is a noun.

\paragraph{Correction 2}
\textbf{Wrong:} The frame is too rigid to bend easily, which means it is very flexible.

\textbf{Correct:} The frame is too rigid to bend easily, which means it is not very flexible.

\textbf{Why:} Rigid and flexible have opposite meanings in physical descriptions.

\paragraph{Correction 3}
\textbf{Wrong:} The material has a rigid.

\textbf{Correct:} The material is rigid.

\textbf{Why:} Rigid is an adjective, not a countable noun.

\subsection{Meaning 2: Strict or difficult to change}
\textbf{Part of speech:} adjective

\textbf{IPA:} /ˈrɪdʒɪd/

\textbf{Pronunciation phrase:} rigid

\textbf{Local audio file:} audios/rigid.mp3

\textbf{Register:} neutral to formal; common in academic, organizational, social, political, and educational contexts

\textbf{Definition:} very strict, fixed, or unwilling to change, even when different circumstances might require flexibility

\textbf{In simpler words:} too fixed or strict

\textbf{Typical pattern:} rigid + rule, system, schedule, structure, approach, hierarchy, or attitude

\subsubsection{Natural examples}
\begin{enumerate}
  \item A rigid schedule can make it difficult for employees to respond to unexpected problems.
  \item The school abandoned its rigid rules and gave teachers greater freedom in the classroom.
  \item Critics argued that the government's rigid approach prevented it from adapting to changing economic conditions.
  \item The organization once had a rigid hierarchy in which junior employees had little influence.
  \item A rigid interpretation of the law may produce unfair results in unusual cases.
  \item The researchers warned against using a rigid classification system for such a complex phenomenon.
  \item His rigid views made productive discussion difficult.
  \item Some companies are replacing rigid working hours with more flexible arrangements.
\end{enumerate}

\subsubsection{Nuance and implication}
\textbf{Central nuance:} This sense transfers the physical idea of something that cannot bend to rules, systems, plans, or attitudes that cannot easily adapt.

\textbf{Usual implication:} The word is often negative because too little flexibility can make a system or person unable to respond well to new circumstances.

\textbf{Compared with flexible:} A flexible rule, system, or approach can be adjusted when circumstances change. A rigid one remains fixed and allows little adjustment.

\subsubsection{Grammar notes}
\textbf{How to use it:} Use it before nouns such as rule, schedule, system, approach, hierarchy, structure, or attitude, or after a linking verb.

\textbf{What can follow it:} It frequently describes policies, institutions, ways of thinking, organizational structures, and procedures.

\textbf{Position in a sentence:} Too rigid and overly rigid are common when criticizing a lack of flexibility.

\subsubsection{Useful patterns}
\textbf{Pattern:} a rigid system or structure

\textbf{Use:} Use this for an organization or arrangement that allows little change.

\textbf{Example:} The rigid administrative system made reform extremely slow.

\textbf{Pattern:} a rigid schedule

\textbf{Use:} Use this for a timetable that allows little or no adjustment.

\textbf{Example:} Parents with young children may struggle with a rigid work schedule.

\textbf{Pattern:} a rigid approach to + noun or -ing form

\textbf{Use:} Use this for a method that does not adapt easily to different situations.

\textbf{Example:} A rigid approach to teaching may not meet the needs of every student.

\textbf{Pattern:} too or overly rigid

\textbf{Use:} Use this when strictness or lack of adaptability creates problems.

\textbf{Example:} The regulations were criticized for being overly rigid.

\textbf{Pattern:} rigid views or beliefs

\textbf{Use:} Use this for opinions that a person is unwilling to reconsider.

\textbf{Example:} Rigid beliefs can make compromise difficult.

\subsubsection{Common wrong usage}
\paragraph{Correction 1}
\textbf{Wrong:} The company follows rigidly rules.

\textbf{Correct:} The company follows rigid rules.

\textbf{Why:} Use the adjective rigid before a noun such as rules.

\paragraph{Correction 2}
\textbf{Wrong:} The schedule is rigid to changes.

\textbf{Correct:} The schedule is rigid and difficult to change.

\textbf{Why:} Rigid is not normally followed by to changes in this meaning.

\paragraph{Correction 3}
\textbf{Wrong:} The manager was rigidly about working hours.

\textbf{Correct:} The manager was rigid about working hours.

\textbf{Why:} Use the adjective rigid after was. Rigidly is an adverb.

\paragraph{Correction 4}
\textbf{Wrong:} A rigid policy is always better because everyone follows the same rule.

\textbf{Correct:} A rigid policy applies fixed rules but may be unsuitable when circumstances differ.

\textbf{Why:} Rigid does not mean fair or effective; it emphasizes lack of flexibility.

\section{Word family}
\sectionrule
\subsection{rigidity}
\textbf{Part of speech:} uncountable noun

\textbf{IPA:} /rɪˈdʒɪdəti/

\textbf{Definition:} the quality of being physically stiff or of being too fixed and unwilling or unable to change

\textbf{Relationship to the headword:} This noun names the physical stiffness or lack of flexibility described by rigid.

\textbf{Grammar pattern:} the rigidity of + object or system

\textbf{Usage note:} It is normally uncountable.

\subsubsection{Examples}
\begin{enumerate}
  \item The rigidity of the structure helps it resist strong winds.
  \item Critics blamed the system's rigidity for its inability to respond quickly to the crisis.
  \item Excessive organizational rigidity can discourage innovation.
\end{enumerate}

\subsection{rigidly}
\textbf{Part of speech:} adverb

\textbf{IPA:} /ˈrɪdʒɪdli/

\textbf{Definition:} in a strict, fixed, or inflexible way

\textbf{Relationship to the headword:} This adverb describes an action performed in a rigid or inflexible way.

\textbf{Grammar pattern:} rigidly + enforce, follow, apply, control, or define

\subsubsection{Examples}
\begin{enumerate}
  \item The organization rigidly enforced its traditional procedures.
  \item The categories should not be applied too rigidly because some cases overlap.
\end{enumerate}

\section{Common collocations}
\sectionrule
\subsection{rigid structure}
\textbf{Applies to:} Stiff and unable to bend

\textbf{Meaning:} a structure that strongly resists bending or changing shape

\textbf{Example:} The bridge requires a rigid structure that can withstand strong winds.

\subsection{rigid material}
\textbf{Applies to:} Stiff and unable to bend

\textbf{Meaning:} a material that does not bend easily

\textbf{Example:} The device is made from a rigid but lightweight material.

\subsection{rigid frame}
\textbf{Applies to:} Stiff and unable to bend

\textbf{Meaning:} a stiff supporting structure

\textbf{Example:} A rigid frame supports the outer shell of the vehicle.

\subsection{rigid surface}
\textbf{Applies to:} Stiff and unable to bend

\textbf{Meaning:} a surface that does not easily bend or deform

\textbf{Example:} The instrument must be placed on a rigid surface.

\subsection{rigid rules}
\textbf{Applies to:} Strict or difficult to change

\textbf{Meaning:} rules that allow very little flexibility

\textbf{Example:} Rigid rules prevented employees from responding creatively to unusual situations.

\subsection{rigid schedule}
\textbf{Applies to:} Strict or difficult to change

\textbf{Meaning:} a timetable that is difficult to adjust

\textbf{Example:} The rigid schedule left no time for unexpected delays.

\subsection{rigid system}
\textbf{Applies to:} Strict or difficult to change

\textbf{Meaning:} a system whose rules or structure are difficult to alter

\textbf{Example:} The old rigid system was gradually replaced by a more flexible one.

\subsection{rigid hierarchy}
\textbf{Applies to:} Strict or difficult to change

\textbf{Meaning:} a strongly fixed ranking system with clearly separated levels

\textbf{Example:} A rigid hierarchy separated senior managers from ordinary employees.

\subsection{rigid approach}
\textbf{Applies to:} Strict or difficult to change

\textbf{Meaning:} a method that allows little adjustment

\textbf{Pattern:} a rigid approach to + noun or -ing form

\textbf{Example:} A rigid approach to language teaching may ignore differences among learners.

\subsection{rigid thinking}
\textbf{Applies to:} Strict or difficult to change

\textbf{Meaning:} a way of thinking that resists new ideas or perspectives

\textbf{Example:} Rigid thinking can make it difficult to consider alternative explanations.

\subsection{too rigid}
\textbf{Applies to:} Strict or difficult to change

\textbf{Meaning:} so inflexible that it causes difficulties

\textbf{Pattern:} be too rigid to + verb

\textbf{Example:} The classification system is too rigid to describe every case accurately.

\subsection{institutional rigidity}
\textbf{Applies to:} Strict or difficult to change

\textbf{Meaning:} the inability of an institution to adapt easily

\textbf{Example:} Institutional rigidity slowed the implementation of the reforms.

\textbf{Usage note:} Rigidity is the noun form.

\section{Synonyms and antonyms}
\sectionrule
\subsection{Close synonyms}
\subsubsection{stiff}
\textbf{Part of speech:} adjective

\textbf{Applies to:} Stiff and unable to bend

\textbf{Shared meaning:} Both can describe something that does not bend easily.

\textbf{Difference:} Stiff is the more common everyday word for something difficult to bend or move. Rigid is more formal and often suggests stronger resistance to changing shape.

\textbf{Example:} The cardboard becomes stiff when several layers are glued together.

\subsubsection{inflexible}
\textbf{Part of speech:} adjective

\textbf{Applies to:} all senses

\textbf{Shared meaning:} Both can describe something that does not adapt or bend easily.

\textbf{Difference:} Inflexible is especially common for rules, plans, attitudes, and people that cannot or will not change. Rigid can describe both this figurative lack of flexibility and literal physical stiffness.

\textbf{Example:} The company's inflexible policies made it difficult for employees with family responsibilities.

\subsubsection{strict}
\textbf{Part of speech:} adjective

\textbf{Applies to:} Strict or difficult to change

\textbf{Shared meaning:} Both can describe rules or systems that place strong limits on behavior.

\textbf{Difference:} Strict means demanding that rules be followed carefully. Rigid goes further by suggesting that the rules or approach allow too little adjustment when circumstances change.

\textbf{Example:} The laboratory follows strict safety procedures.

\subsubsection{fixed}
\textbf{Part of speech:} adjective

\textbf{Applies to:} Strict or difficult to change

\textbf{Shared meaning:} Both can describe something that does not change easily.

\textbf{Difference:} Fixed means established and not changing. Rigid often adds a negative idea that the lack of change creates difficulties or prevents adaptation.

\textbf{Example:} The course follows a fixed timetable throughout the semester.

\subsection{Direct or useful antonyms}
\subsubsection{flexible}
\textbf{Part of speech:} adjective

\textbf{Applies to:} all senses

\textbf{Opposition:} Flexible describes something that can bend physically or adapt to changing circumstances, while rigid describes something that strongly resists bending or change.

\textbf{Example:} The flexible material can bend without breaking.

\subsubsection{adaptable}
\textbf{Part of speech:} adjective

\textbf{Applies to:} Strict or difficult to change

\textbf{Opposition:} An adaptable person, system, or approach can change to suit new circumstances, while a rigid one changes with difficulty.

\textbf{Example:} Successful organizations need adaptable systems that can respond to unexpected conditions.

\section{Words learners may confuse}
\sectionrule
\subsection{rigorous}
\textbf{Part of speech:} adjective

\textbf{Why learners confuse them:} Both are formal adjectives beginning with rig- and can describe demanding systems or procedures, but their meanings are different.

\textbf{Key difference:} Rigid means physically stiff or unwilling or unable to change. Rigorous means extremely careful, thorough, or demanding.

\textbf{rigid:} The rigid schedule allowed almost no changes.

\textbf{rigorous:} The researchers used a rigorous method to test the hypothesis.

\subsection{strict}
\textbf{Part of speech:} adjective

\textbf{Why learners confuse them:} Both can describe rules and systems that strongly limit behavior.

\textbf{Key difference:} Strict emphasizes strong rules and careful enforcement. Rigid emphasizes lack of flexibility and often suggests that adjustment is difficult or impossible.

\textbf{rigid:} The rigid policy could not be adapted to unusual cases.

\textbf{strict:} The laboratory has strict rules about handling dangerous chemicals.

\section{Etymology}
\sectionrule
\textbf{Origin:} Rigid comes from Latin rigidus, meaning stiff or hard.

\textbf{Memory link:} The original idea of physical stiffness helps explain the figurative meaning: a rigid rule or attitude cannot easily bend or adapt.

\section{Practice exercises}
\sectionrule
\subsection{Word family choice}
Choose the best member of the word family for each blank.

\begin{enumerate}
  \item The \_\_\_\_\_ of the old administrative system made rapid reform difficult.
  \item The categories should not be applied too \_\_\_\_\_.
  \item The container must be sufficiently \_\_\_\_\_ to protect the equipment.
\end{enumerate}

\subsection{Correct or incorrect?}
Decide whether each sentence is correct. Correct the inaccurate ones.

\begin{enumerate}
  \item The rigid steel frame helps the structure maintain its shape.
  \item The company follows a very rigidity schedule.
  \item A rigid approach may become ineffective when conditions change.
  \item The researchers used a rigid method, meaning that it was extremely careful and thorough.
  \item The organization rigidly enforced the same procedure in every case.
\end{enumerate}

\subsection{Collocation practice}
Complete each sentence with a natural collocation.

\begin{enumerate}
  \item The company replaced its rigid \_\_\_\_\_ with a more flexible timetable.
  \item The machine requires a \_\_\_\_\_ frame to keep all of its components properly aligned.
  \item A rigid approach \_\_\_\_\_ teaching may not work equally well for every student.
  \item Excessive institutional \_\_\_\_\_ can prevent organizations from adapting to change.
\end{enumerate}

\subsection{Complete answer key}
\subsubsection{Word family choice}
\begin{enumerate}
  \item \textbf{rigidity.} The rigidity of the old administrative system made rapid reform difficult. \emph{Use the noun rigidity for the quality of being inflexible.}
  \item \textbf{rigidly.} The categories should not be applied too rigidly. \emph{Use the adverb rigidly to describe how the categories are applied.}
  \item \textbf{rigid.} The container must be sufficiently rigid to protect the equipment. \emph{Use the adjective rigid after be.}
\end{enumerate}

\subsubsection{Correct or incorrect?}
\begin{enumerate}
  \item \textbf{Correct} \emph{Rigid correctly describes a frame that strongly resists bending.}
  \item \textbf{Incorrect}. The company follows a very rigid schedule. \emph{Use the adjective rigid before the noun schedule.}
  \item \textbf{Correct} \emph{Rigid appropriately describes an approach that does not adapt easily.}
  \item \textbf{Incorrect}. The researchers used a rigorous method, meaning that it was extremely careful and thorough. \emph{Rigorous means careful and thorough. Rigid means stiff or inflexible.}
  \item \textbf{Correct} \emph{Rigidly correctly describes inflexible enforcement.}
\end{enumerate}

\subsubsection{Collocation practice}
\begin{enumerate}
  \item \textbf{schedule.} The company replaced its rigid schedule with a more flexible timetable. \emph{A rigid schedule is one that allows little adjustment.}
  \item \textbf{rigid.} The machine requires a rigid frame to keep all of its components properly aligned. \emph{A rigid frame strongly resists bending or changing shape.}
  \item \textbf{to.} A rigid approach to teaching may not work equally well for every student. \emph{Use an approach to followed by a noun or -ing form.}
  \item \textbf{rigidity.} Excessive institutional rigidity can prevent organizations from adapting to change. \emph{Institutional rigidity means lack of flexibility within an institution.}
\end{enumerate}

\vfill
\begin{center}
\small\color{Muted} Created and compiled by \textbf{Amir Kooshky}\\
\href{https://t.me/KooshkyTOEFL}{Telegram Channel}\quad\textbullet\quad
\href{https://www.instagram.com/kooshkytoefl}{Instagram}\quad\textbullet\quad
\href{https://t.me/KooshkyTOEFL\_pv}{Personal Telegram}
\end{center}
\end{document}
